\begin{proposition}\label{sga3:XXI:3.2.6}
If \(\Delta\) is a system of simple roots, then
\(\mathscr{P}(\Delta)\) is a system of positive roots. If \(\Delta\)
is a system of simple roots and \(R_+\) is a system of positive roots,
then
\[
 \Delta\subset R_+
 \quad\Longleftrightarrow\quad
 R_+=\mathscr{P}(\Delta).
\]
\end{proposition}

\begin{proof}
The first assertion is immediate. If \(\Delta\subset R_+\), then
\(\mathscr{P}(\Delta)\subset R_+\) by (P 3), and therefore
\(\mathscr{P}(\Delta)=R_+\) by
Remark~\ref{sga3:XXI:3.2.3}. The rest is trivial.
\end{proof}

\begin{remark}\label{sga3:XXI:3.2.7}
Systems of positive roots exist: endow \(\VR(R)\) with the structure
of a totally ordered vector space. The set of roots that are
\(\geq0\) for this order is a system of positive roots.
\end{remark}

\begin{theorem}\label{sga3:XXI:3.2.8}
Let \(R_+\) be a system of positive roots. There exists a unique
system of simple roots \(\mathscr{S}(R_+)\) such that
\[
 \mathscr{S}(R_+)\subset R_+,
 \qquad\text{i.e.}\qquad
 R_+=\mathscr{P}\bigl(\mathscr{S}(R_+)\bigr).
\]
\end{theorem}

Uniqueness follows immediately from:

\begin{lemma}\label{sga3:XXI:3.2.9}
Let \(\Delta\) be a system of simple roots. An element
\(\alpha\in\mathscr{P}(\Delta)\) belongs to \(\Delta\) if and only if
\(\alpha\) is not the sum of two elements of
\(\mathscr{P}(\Delta)\).
\end{lemma}

\begin{proof}
This follows immediately from the definitions and
Proposition~\ref{sga3:XXI:3.1.2}.
\end{proof}

We now prove the existence of \(\mathscr{S}(R_+)\). Consider the set
of subsets \(T\) of \(R_+\) such that
\[
 T=\operatorname{ind}(T)
 \qquad\text{and}\qquad
 \Qp\cdot T\supset R_+.
\]
This set is nonempty, since it contains
\(\operatorname{ind}(R_+)\). Let \(\Delta\) be a minimal element of
this set under inclusion. We show that \(\Delta\) is a system of
simple roots; by Proposition~\ref{sga3:XXI:3.1.5}\SGARefLink{sga3:XXI:3.1.5:item:i}{(i)}, it is enough to
show that \(\Delta\) is a linearly independent subset of
\(M\otimes\mathbb Q\).

\begin{lemma}\label{sga3:XXI:3.2.10}
If \(\alpha,\beta\in\Delta\) and
\[
 q\alpha+q'\beta\in R,
 \qquad q,q'\in\mathbb Q,
\]
then \(qq'\geq0\).
\end{lemma}

\begin{proof}
If \(qq'<0\), then, after interchanging \(\alpha\) and \(\beta\) if
necessary, we may write
\[
 a\alpha-b\beta\in R_+,
 \qquad a,b\in\mathbb Q^\ast_+.
\]
By the defining property of \(\Delta\), there is a relation
\[
 a\alpha-b\beta
   =\sum_{\gamma\in\Delta}c(\gamma)\gamma,
 \qquad c(\gamma)\in\Qp.
\]
If \(a\leq c(\alpha)\), this becomes
\[
 -\beta
 =\frac{c(\alpha)-a}{b}\alpha
  +\sum_{\gamma\neq\alpha}\frac{c(\gamma)}{b}\gamma.
\]
Thus \(-\beta\in(\Qp\cdot\Delta)\cap R\), which is contained in
\(R_+\) by (P 3). This gives
\(\beta\in R_+\cap(-R_+)\), contradicting (P 1).

If instead \(a>c(\alpha)\), then
\[
 \bigl(a-c(\alpha)\bigr)\alpha
 =b\beta+\sum_{\gamma\neq\alpha}c(\gamma)\gamma,
\]
which shows that
\[
 \Delta\subset\Qp\cdot\bigl(\Delta\setminus\{\alpha\}\bigr),
\]
contrary to the minimality of \(\Delta\).
\end{proof}
